% No preamble. Prefix mgs-. Flat mean-zero-gauge spectral estimates.
\section{Both flat mean-zero-gauge factors have a physical-scale gap}
\label{mgs-section}

This section proves spectral estimates for the explicit flat matrices
of the corrected all-site gauge construction.
The exact nonlinear determinant identity and its continuation through
the reducible flat background are separate inputs.
No derivative estimate for a varying-background heat trace is inferred
from these flat bounds.

\subsection{Spaces and flat operators}
\label{mgs-spaces-section}

Work in dimension \(d=4\), on a cubic current torus of side \(KN\),
with \(K=\ell^k\) odd, final centers \(Ky\), and the disjoint centered
cubes \(B_y\) of side \(K\).
All scalar and bond norms are counting norms.
Let \(D\) be the full-site scalar gradient, \(d\) the curl, and
\(\Delta_s=D^TD\).
Let \(Q_0\) be the normalized scalar block average, and put
\begin{equation}
 P=I-Q_0^T(Q_0Q_0^T)^{-1}Q_0,\qquad \Pi=I-P .
 \label{mgs-projections}
\end{equation}
Thus \(P\) projects onto block-mean-zero scalars, and \(\Pi\)
onto block-constant scalars.
Choose a block-diagonal orthonormal frame \(V\) for \(\ker Q_0\)
and define
\begin{equation}
                       G_B=V^T\Delta_sV.
 \label{mgs-ghost}
\end{equation}
At \(K=1\) this scalar space is empty; its trace is zero and
its determinant is one.

The actual flat corrected composite differential is the straight
angle average \(A=A_K\). It obeys
\begin{equation}
 AD=D_cQ_0,\qquad \|A\|_{2\to2}\le K^{-1}.
 \label{mgs-average}
\end{equation}
Here \(D_c\) is the coarse scalar gradient.
The estimate is for this corrected differential; it is not
asserted for the full ordinary unprojected constraint.

Let \(L=q_M\) and \(q=q_{MK}\) be the actual flat classical
Hessians. The established homogeneous inputs are
\begin{equation}
 L\ge d^Td,\qquad LD=0,\qquad 0\le q\le R_qI,
 \label{mgs-classical-inputs}
\end{equation}
uniformly in old depth \(M\), endpoint \(K\), and volume.
Choose a fixed \(a\ge R_q+1\), and set
\begin{equation}
 Q_B=L-A^TqA+DPD^T+aA^TA.
 \label{mgs-vector}
\end{equation}
Its energy satisfies
\begin{equation}
 \langle x,Q_Bx\rangle
 \ge \|dx\|^2+\|PD^Tx\|^2+\|Ax\|^2.
 \label{mgs-energy-lower}
\end{equation}
Matching these matrices to the nonlinear bridge uses only its
flat identifications \(B(1)=Q_0\) and \(\widehat C(1)=A_K\).

\subsection{Scalar Poincare and a block-mean harmonic lift}
\label{mgs-scalar-section}

If \(Q_0\psi=0\), the Neumann cube Poincare inequality, summed
over the blocks, gives
\begin{equation}
 \|\psi\|^2\le C_PK^2\|D\psi\|^2,\qquad C_P=1.
 \label{mgs-poincare}
\end{equation}
Indeed the first nonzero cube eigenvalue is
\(4\sin^2(\pi/(2K))\ge4/K^2\), and cube-internal edges
are a subset of all edges.
The case \(K=1\) is vacuous.
In particular,
\begin{equation}
                         G_B\ge C_P^{-1}K^{-2}I.
 \label{mgs-ghost-gap}
\end{equation}

For a coarse scalar \(b\), let \(H_b\) minimize
\(\|D\phi\|^2\) subject to \(Q_0\phi=b\).
The constraint is onto and fixes the constant nullspace, so the
finite minimizer exists uniquely. Its Euler equation is
\(P\Delta_sH_b=0\).
Moreover,
\begin{equation}
 \|DH_b\|^2\le C_{\rm lift}K^2\|D_cb\|^2,\qquad
 C_{\rm lift}=\frac{\pi^2}{4}\left(\frac\pi2\right)^8.
 \label{mgs-harmonic-upper}
\end{equation}
To verify the normalization, lift a coarse frequency
\(p\in[-\pi,\pi]^4\) to the central fine frequency \(p/K\).
The scalar average has multiplier, up to a unit phase,
\[
 a_0(p)=\prod_{\mu=1}^4
      \frac{\sin(p_\mu/2)}{K\sin(p_\mu/(2K))},
 \qquad |a_0(p)|\ge(2/\pi)^4.
\]
Divide the lifted amplitude by \(a_0\).
The fine difference obeys
\[
 |e^{ip_\mu/K}-1|
       \le\frac{\pi}{2K}|e^{ip_\mu}-1|.
\]
Counting Parseval multiplies the fine squared amplitude by
\(K^4\), relative to its coarse counterpart, proving
\eqref{mgs-harmonic-upper}.
The zero mode is constant and has zero energy.
All selected frequencies are legitimate on the \(KN\)-periodic
torus. For real data one may take the real part of a complex
feasible lift; this preserves the constraint and does not
increase energy, including at coarse Nyquist frequencies.

For arbitrary \(\phi\), put \(b=Q_0\phi\) and
\(\psi=\phi-H_b\). Then
\[
 Q_0\psi=0,\qquad
 P\Delta_s\psi=P\Delta_s\phi=:g,\qquad
 \|D\psi\|^2=\langle\psi,g\rangle.
\]
Equation \eqref{mgs-poincare} therefore gives
\(\|D\psi\|\le\sqrt{C_P}K\|g\|\), and hence
\begin{equation}
 \|D\phi\|^2
 \le 2C_PK^2\|P\Delta_s\phi\|^2
           +2C_{\rm lift}K^2\|D_cQ_0\phi\|^2.
 \label{mgs-scalar-control}
\end{equation}

\subsection{The full vector gap}
\label{mgs-gap-section}

\begin{theorem}[Uniform flat mean-zero-gauge coercivity]
\label{mgs-gap-theorem}
There is a constant \(C_{\rm gap}\), independent of \(M,K,N\),
such that
\begin{equation}
 \|x\|^2\le C_{\rm gap}K^2
       \{\|dx\|^2+\|PD^Tx\|^2+\|Ax\|^2\},
 \qquad
 Q_B\ge C_{\rm gap}^{-1}K^{-2}I .
 \label{mgs-vector-gap}
\end{equation}
\end{theorem}

\begin{proof}
Use the orthogonal periodic Hodge decomposition
\(x=x_T+D\phi\), where \(D^Tx_T=0\).
Harmonic one-forms are included in \(x_T\).
The previously proved transverse alias estimate is
\begin{equation}
 \|x_T\|^2\le K^2
       \{A_0\|P_cAx_T\|^2+B_0\|dx_T\|^2\},
 \quad
 A_0=2c_\omega^{-2},\quad
 B_0=(1+2c_\omega^{-2})/4,\quad c_\omega=(2/\pi)^5.
 \label{mgs-transverse}
\end{equation}
It is a Maxwell/averaging inequality, includes zero momentum,
and has no old-depth dependence.

By \eqref{mgs-average},
\[
 P_cAx_T=P_cAx,\qquad
 \|D_cQ_0\phi\|^2
 \le2\|Ax\|^2+2K^{-2}\|x_T\|^2.
\]
Also \(P\Delta_s\phi=PD^Tx\) and \(dx_T=dx\).
Insert these identities in \eqref{mgs-scalar-control},
then add \(\|x_T\|^2\) and apply \eqref{mgs-transverse}.
A valid constant is
\begin{equation}
 C_{\rm gap}=
 \max\{(1+4C_{\rm lift})B_0,\ 2C_P,\
               (1+4C_{\rm lift})A_0+4C_{\rm lift}\}.
 \label{mgs-gap-constant}
\end{equation}
This proves the first inequality in \eqref{mgs-vector-gap};
\eqref{mgs-energy-lower} proves the second.
\end{proof}

No point-pinned scalar gap or inverse of a globally based scalar
matrix occurs in this argument.

\subsection{A local tent lift controls the full bond Laplacian}
\label{mgs-tent-section}

\begin{lemma}[Block-constant scalar lifting]
\label{mgs-tent-lemma}
There is a block-local map \(E:\operatorname{ran}\Pi\to
\{\text{all scalar sites}\}\) such that
\begin{equation}
 \Pi E=I,\qquad \|E\|\le C_0,\qquad
 \|DE\|\le C_1/K,\qquad
 C_0=2^{d/2},\quad C_1^2=80d\,2^{d-1}.
 \label{mgs-tent-bounds}
\end{equation}
These bounds include all periodic multiplicities.
\end{lemma}

\begin{proof}
For \(K=2h+1\ge3\), put
\[
 f(r)=\frac K{h+1}\left(1-\frac{|r|}{h+1}\right),
 \qquad -h\le r\le h.
\]
Its average is one, its mean square is at most two, and its
endpoint value and adjacent differences are at most \(4/K\).
If \(e\) equals \(c_y\) on \(B_y\), define
\[
 (Ee)(Ky+r)=c_y\prod_{\mu=1}^d f(r_\mu).
\]
The block mean is \(c_y\), and tensorization gives the stated
norm bound. For a fixed direction, internal edges contribute
at most
\(16\,2^{d-1}K^{d-2}\sum_y|c_y|^2\)
to the squared gradient norm.
Interfaces contribute at most
\(64\,2^{d-1}K^{d-3}\sum_y|c_y|^2\), because
\(\sum_y|c_{y+e_\mu}-c_y|^2\le4\sum_y|c_y|^2\).
Divide by
\(\|e\|^2=K^d\sum_y|c_y|^2\) and sum directions.
The same estimates count the repeated periodic interfaces when
\(N=1\) or \(2\).
For \(K=1\), take \(E=I\); the conservative constants still apply.
\end{proof}

For a bond field \(x\) and \(e\in\operatorname{ran}\Pi\),
the equality \(Ee-e\in\operatorname{ran}P\) gives
\[
 \langle\Pi D^Tx,e\rangle
  =\langle x,DEe\rangle-\langle PD^Tx,Ee-e\rangle.
\]
Thus
\begin{equation}
 \|\Pi D^Tx\|
 \le\frac{C_1}{K}\|x\|+(C_0+1)\|PD^Tx\|.
 \label{mgs-duality}
\end{equation}
This step uses only a local lifting and duality.

The full bond Laplacian satisfies
\(\langle x,\Delta_1x\rangle=\|dx\|^2+\|D^Tx\|^2\).
Combining \eqref{mgs-duality},
\eqref{mgs-energy-lower}, and \eqref{mgs-vector-gap}
proves
\begin{equation}
 \begin{split}
 \Delta_1&\le C_\Delta Q_B,\qquad
 C_\Delta=1+2(C_0+1)^2+2C_1^2C_{\rm gap},\\
 Q_B&\ge \frac1{2C_\Delta}\Delta_1
                         +\frac1{2C_{\rm gap}K^2}I.
 \end{split}
 \label{mgs-laplacian-domination}
\end{equation}

\subsection{Uniform heat density and inverse moments for both factors}
\label{mgs-heat-section}

\begin{corollary}[Flat local-gauge heat bounds]
\label{mgs-heat-corollary}
For either \(O=Q_B\) or \(O=G_B\),
\begin{equation}
 N^{-4}\operatorname{Tr}e^{-uK^2O}
 \le C\left\{\frac{K^4}{(1+K^2u)^2}+N^{-4}\right\}e^{-cu},
 \qquad u>0,
 \label{mgs-heat}
\end{equation}
uniformly in \(M,K,N\).
The trace for \(G_B\) is in its block-local
\(\ker Q_0\) frame.
In the infinite flat structural system the \(N^{-4}\) term is
omitted. Both factors obey the inverse-moment bounds and the
weighted final-time exponential tail of
Theorem~\ref{pvs-main}, with \(O\) in place of its physical \(Q\).
\end{corollary}

\begin{proof}
For \(Q_B\), apply eigenvalue min--max to
\eqref{mgs-laplacian-domination}.
The comparison is the full scalar lattice heat trace, with a
fixed rescaling of time and multiplicity \(4n_{\mathfrak g}\),
times \(e^{-cs/K^2}\).

For \(G_B\), \eqref{mgs-ghost-gap} gives
\[
 G_B\ge\tfrac12G_B+(2C_PK^2)^{-1}I.
\]
The eigenvalues of the compression \(V^T\Delta_sV\)
dominate the correspondingly indexed full scalar eigenvalues,
by the variational principle.
Its heat trace is therefore at most the full scalar heat trace.
The gap supplies the same exponential factor.

The scalar heat estimate and the substitution \(s=K^2u\) are
exactly those used in Theorem~\ref{pvs-main}.
Final-block translation invariance identifies the normalized
total with one block trace.
Coarse Bloch integration proves the infinite version directly.
Mellin integration gives the stated inverse powers and tails.
\end{proof}

The basis-free scalar completion used by the analytic determinant
bridge is
\begin{equation}
 \widehat G_B=P\Delta_sP+K^{-2}(I-P).
 \label{mgs-basis-free-scalar}
\end{equation}
At flatness its two subspaces reduce it. Its heat density in final
time is exactly
\[
 N^{-4}\operatorname{Tr}e^{-uK^2\widehat G_B}
 =N^{-4}\operatorname{Tr}e^{-uK^2G_B}
                        +n_{\mathfrak g}e^{-u}.
\]
This extra term satisfies \eqref{mgs-heat}, after decreasing \(c\):
for \(u\le1\),
\(K^4/(1+K^2u)^2\ge1/4\), while for \(u\ge1\) it is at least
\(1/(4u^2)\), and \(u^2e^{-(1-c)u}\) is bounded when \(c<1\).
The normalized inverse moment adds exactly \(n_{\mathfrak g}\),
which is absorbed in every case of \eqref{pvs-moments}.
Thus the actual basis-free pair \(\widehat G_B,Q_B\) has the
same heat, inverse-moment, and final-time tail bounds.

These estimates concern the explicit flat local-gauge factors.
They do not establish uniform relative-form derivatives under an
affine curvature source, or identify separate massless factors
with independent continuum references.
The exact nonlinear determinant bridge and the remaining
large-\(K\) curvature comparison must retain their own
hypotheses and proofs.

